\section{A coherent book needs complete collections and ordered identities}
\label{sec:collection-state-requirements}

The return equations can be arithmetically correct while the proposed sequence
is impossible. A sale may precede the creation of its recovered inventory lot;
the same lot may be sold twice; or a balanced book may contain negative cash.
Each failure concerns information that a scalar equality does not express. The
finite collection and state references extend the preceding requirements to
those relationships, retaining a distinction between checking a supplied state
and applying an accounting transition.

In the retained example, settlement closes the return component and creates a
recovered lot. A later sale consumes that lot. The lot's carrying amount also
contributes to inventory in the shared legal book. Three different checks follow:
the relevant records must be present and identifiable, their sequence must make
the lot available when it is sold, and their amounts must reconcile with the
book. Passing any one of these checks cannot establish the other two.

\subsection{Which members are actually known}

\code{MethodCollections/1} declares typed collections and selected constraints
over their fields. Its first operations are uniqueness, exact-one reference
resolution, nondecreasing order, and a guarded requirement evaluated for each
row. Global inputs and row inputs have explicit mappings to the preserved
guarded evaluator. Type, context and role must agree; a field name does not
by itself establish economic equivalence.

A supplied collection is either complete with its rows, or partial with its
observed rows and a reason for incomplete coverage. Missing collections are
partial with no observed rows. Partial rows retain their relative supplied
order, while additional members may be absent between them. An unknown field
has a separate tagged reason. Thus an empty complete list says that there are
no members in the declared scope; an empty partial list says that none has yet
been supplied. These are different assertions, and neither proves that the
scope describes an actual company's complete record.

Known duplicates contradict uniqueness even when another key is unknown. When
no duplicate is observed, an unknown key or incomplete membership prevents a
uniqueness conclusion. A reference must resolve to exactly one matching target
of the same type and context. Two known matches already contradict that
requirement. No match is a contradiction only when the target collection is
complete and its keys are known. One known match remains unresolved if another
unknown or unobserved target could create a second match. Target-wide uniqueness
is a separate condition when needed.

Ordering checks preserve the caller's sequence. If two known dates are reversed,
an unknown date between them cannot repair the reversal. If the known subsequence
is ordered but another date or member is unknown, the full order remains
unresolved. Equal dates retain their original order rather than being sorted
by a hidden tie-breaker. Period predicates separately check the retained
half-open interval: an event on the start date is included, while an event on
the exclusive end date is outside the period.

Per-row checks retain all original guarded findings and indexed input mappings.
A known invalid row remains visible despite another unknown row or incomplete
membership. An inactive condition still undergoes type and structure checks,
and its separately required facts remain required. A complete empty collection
vacuously satisfies a universal row condition; a partial empty collection leaves
it unresolved. These rules prevent missing records from being mistaken for a
successful review of every record.

\subsection{A sale needs an earlier live lot}

Snapshot membership cannot show that a lot existed before a sale. The retained
comparison makes this concrete by giving settlement and sale the same date and
placing the sale first. Date ordering then passes, but the native method rejects
the sale because the lot does not yet exist. A second sale and a second complete
settlement expose similar limitations. These failures motivate an explicitly
ordered identity check rather than a stronger interpretation of the snapshot.

\code{MethodIdentityLifecycle/1} declares named identity pools, their initial
live and retired members, and a closed action mapping. Each action contains an
ordered list of \code{create}, \code{require} or \code{consume} operations.
Create requires an identity absent from both live and retired memory. Require
checks that an identity is live without changing it. Consume moves a live
identity into retired memory. Retaining retired identities prevents an already
used identifier from being silently created again.

For the conditional return example, settlement consumes the live return-component
identity and creates a recovered-lot identity. Sale consumes that lot. Consuming
the component makes a second complete settlement fail; consuming the lot makes
a second whole-lot sale fail. The supplied order also makes a sale before creation
fail, including when both events have the same date. The initial empty retired
pools in this example are explicit assumptions of its synthetic opening. They
are not evidence about an actual company's lifetime event history.

Each event's operations run against a private copy of the identity pools. Only
when every operation succeeds does the diagnostic trace adopt that event's
result. If a later operation fails or needs an unknown key, no partial identity
change from that event survives. Earlier successful steps remain visible with
their indices and after-states. This is local identity atomicity: no money is
posted, no journal is authorized and no durable transaction is committed.

Unknown initial identities or incomplete initial pools stop the trace before its
first event. An incomplete event sequence also stops before any observed event,
because an omitted predecessor might create or consume a needed resource. With
a complete sequence, an unknown action or operation key stops at that event.
Later state is not inferred. All observed rows still undergo structural checks,
and independent snapshot findings remain available. An otherwise complete
identity trace can therefore coexist with an unresolved snapshot condition;
the combined result does not report success.

The four retained failures---sale before creation on an equal date, a missing
sale lot, repeated whole-lot sale and repeated complete settlement---are all
rejected by this explicit identity sequence. The result establishes only those
declared resource relationships. Prices, cash availability, tax consequences
and accounting entries still require their own conditions.

\subsection{Signed totals and the limits of balance}

The shared legal book uses signed account balances. In the retained convention,
assets normally carry positive coefficients and liabilities and equity negative
coefficients. For the declared complete account registry $\mathcal A$, the
selected balance condition is
\[
  \sum_{a\in\mathcal A} b_a=0.
\]
The registry matters as much as the sum. Omitting an account or adding an
unclassified zero-balance account is a structural mismatch even if the displayed
coefficients still add to zero. The retained declaration therefore checks unique
account keys and reference membership in both directions against the pinned
native account list. Account-classification coverage is checked separately
against the implementation's classification lists.

Let $\mathcal L$ be the declared complete recovered-lot collection, $c_\ell$ its
carrying amount for lot $\ell$, and $I_{\mathrm{other}}$ the separately represented
inventory. Selected allocation conditions include
\begin{align}
 b_{\mathrm{inventory}} &= I_{\mathrm{other}}+
                          \sum_{\ell\in\mathcal L}c_\ell, \\
 b_{\mathrm{refund\ liability}} &= -L_{\mathrm{refund}}, \\
 b_{\mathrm{recovery\ asset}} &= A_{\mathrm{recovery}}.
\end{align}
The minus sign in the refund equation follows the signed-book convention; it
is not an absolute-value normalization. The retained lot tax method separately
requires each lot's tax basis to equal its carrying amount. Other tax methods
need different, explicitly selected conditions.

\code{MethodCollectionTotals/1} binds an assessment input either to a declared
global value or to the sum of one declared rational collection field. It retains
the original type, role and contributor indices. Complete, known inputs give an
exact reduced rational sum; an empty complete collection gives zero. A missing
or unknown field, or partial collection membership, leaves the total unknown.
The observed subtotal remains available as a diagnostic but is never substituted
for that unknown total. Even cancellation of all observed amounts cannot remove
an unknown contribution.

Contexts, dimensions and scales must match exactly. There is no inferred currency
conversion, group-to-legal allocation, rounding or balancing entry. Snapshot
findings and the guarded scalar assessment remain visible together, so a correct
sum cannot override a duplicate account or incomplete reference registry.
Completeness remains a supplied declaration whose source support must be checked
before using the result as a claim about a company.

Two counterexamples show the remaining limitation. One adverse fixture has a
negative cash balance and an offsetting capital change, keeping the book balanced
solely to isolate the funding check. Another marks a lot disposed while leaving
its stock amounts unchanged. Both satisfy the selected totals. Neither is an
admissible state under the retained method. The capital offset belongs only to
the diagnostic fixture; the engine does not manufacture financing to repair a
cash shortfall.

\subsection{Distinguish funding failure from an invalid state}

The additional state predicates check nonnegative cash and other inventory,
settled return rights, lot identities and amount domains, disposition memory,
lifetime event identities and tax-adjustment references. A disposed lot must
have zero units, carrying amount and tax basis; a live lot must have positive
units. Its stored identifier must match its collection key and its component
must match the owner. Lifetime event IDs must be unique, and each uniquely named
tax adjustment must refer to an event already in that history.

These conditions reuse guarded scalar predicates, collection constraints and
exact totals. Known settlement and disposition facts are separately required;
unknown is never interpreted as false. A narrow
\code{MethodStateAssessment/1} layer preserves the economic category of an active,
known-false cash condition. It reports \code{funding\_failure} while retaining the
original predicate result. It cannot relabel malformed values, inactive
conditions, unknowns or unrelated failures as funding outcomes.

The complete findings summary orders invalid state, funding failure, unsupported
method scope and unresolved information, retaining every finding. This differs
from the native invariant's first-failure control flow. If cash is negative and
a later lifetime event ID is duplicated, the native function first reports the
cash shortfall; the common assessment reports invalid state and also retains
the funding failure. Another explicit difference concerns disposition memory:
the native function rejects an unknown disposition because it requires a
concrete Boolean, whereas the tagged-input extension reports unresolved. Neither
difference is hidden inside a claim of full native equivalence.

The saved comparisons repair both counterexamples: balanced negative cash is a
funding failure and disposed stock is invalid. Selected adverse cases also cover
identity mismatches, negative or empty live lots, outstanding settled rights and
invalid tax references. The branch disposition record distinguishes these
checked predicates from fixed implementation constants and the still-separate
validation of complete outer and nested record shape.

The resulting chain gives increasingly informative prerequisite assessments:
typed facts, applicable equations, complete collections, ordered identity memory,
signed reconciliations and state-domain findings. Each stage leaves source and
accounting admission, full method eligibility and execution false. The next section connects these prerequisites to a selected return transition:
complete raw-state structure, exact monetary effects, atomic memory updates and
full proposal recomputation. Reporting and annual carry retain their own method
conditions, and source correspondence remains distinct from authority.
